\section{Linux Networking}\label{linux-networking}

Per approfondire meglio il discorso relativo al routing dei pacchetti e
delle interfacce di rete, si vuole elencare ad alto livello cosa succede
all'interno del kernel linux e quali sono gli oggetti fondamentali.

\subsection[Network Device Driver ]{\texorpdfstring{Network Device
Driver \footnote{\label{fn:ref_23}Linux Kernel Networking - Implementation and Theory}
}{Network Device Driver  }}\label{network-device-driver-1}

Un network device opera sempre livello 2 o 3 nello stack di rete, ed e'
rappresentato nel kernel dalla struttura
\passthrough{\lstinline!net\_device!}. Questa struttura dati contiene
varie informazioni:

\begin{itemize}
\tightlist
\item
  \passthrough{\lstinline!IRQ!}: numero identificativo del dispositivo
\item
  \passthrough{\lstinline!MTU!}: la MTU del dispositivo
\item
  \passthrough{\lstinline!MAC!}: MAC address del dispositivo
\item
  \passthrough{\lstinline!Nome!}: \passthrough{\lstinline!eth0!},
  \passthrough{\lstinline!enp0s3!}, \passthrough{\lstinline!wlan0!},
  ecc\ldots{}
\item
  \passthrough{\lstinline!net\_device\_ops!}: contiene i puntatori alle
  funzioni per la gestione del dispositivo (come le primitive start,
  stop, transmit, receive)
\end{itemize}

Un device driver si occupa di:

\begin{itemize}
\tightlist
\item
  \textbf{Ricevere} i pacchetti destinati al local host, inoltrandoli al
  livello di rete, e successivamente al livello di trasporto.
\item
  \textbf{Trasmette} i pacchetti generati dal local host, o fa il
  forward dei pacchetti ricevuti verso un'altra interfaccia.
\end{itemize}

Ogni pacchetto in transito in tutti e due i sensi scatena un lookup nel
\emph{sottosistema di routing}, che decide cosa bisogna fare con quel
pacchetto. La fase di routing si distingue in 5 fasi, detti
\textbf{hook}, dove possiamo registrare dei \textbf{callback} per il
\textbf{sottosistema netfilter}.

Per esempio l'hook \passthrough{\lstinline!NF\_INET\_PRE\_ROUTING!},
registra la chiamata al callback per essere eseguito prima che il lookup
del sottosistema di routing venga effettuato. Il callback ritorna un
verdetto sul filtraggio del pacchetto, decidendo se bisogna accettarlo
(\passthrough{\lstinline!NF\_ACCEPT!}), oppure scartarlo
(\passthrough{\lstinline!NF\_DROP!}).

\begin{quote}
Con \passthrough{\lstinline!nftables!} applichiamo le regole di
filtraggio dei pacchetti in modo da intercettare e manipolare il
sottosistema di routing
\end{quote}

\begin{quote}
Oltre ai callback registrati, altri fattori influenzano il modo in cui
il routing viene effettuato. Per esempio un pacchetto viene scartato
quando il TTL viene diminuito e arriva a 0.
\end{quote}

La struttura \passthrough{\lstinline!sk\_buff!} (Socket Buffer)
definisce come un pacchetto nel kernel Linux viene rappresentato e viene
manipolato unicamante attraverso la SKB API.

La Socket Buffer viene allocata dinamicamente quando il pacchetto viene
ricevuto sul cavo, ed e' la struttura dati utilizzata per manipolare il
pacchetto nel kernel sui vari layer di rete.

All'interno del kernel la struttura dati responsabile per il
sottosistema di routing e' la \passthrough{\lstinline!fib\_table!},
ossia la tabella che raccoglie le rotte per il traffico destinato ad una
certa sottorete.

\subsection[Bridge ]{\texorpdfstring{Bridge
\footnote{\label{fn:ref_24}Understanding Linux Networking Internals}}{Bridge }}\label{bridge-2}

Un bridge e' un dispositivo L2 che capisce il protocollo del livello di
collegamento, effettua l'inoltro dei dati copiandoli frame per frame.
Per farlo deve mantenere nel suo buffer almeno un frame per volta.

Un bridge assegna un'indirizzo al livello di collegamento ad ognuna
delle sue interfacce ed effettua il forward verso di queste se il bridge
non e' il destinatario del pacchetto. Opzionalmente puo' implementare il
livello di rete (Layer 3) per offrire visibilita alle interfacce
collegate e viene di solito utilizzato per unire due LAN diffrenti, o
per creare una LAN.

Un bridge nel kernel Linux e' un device virtuale, finché non viene
collegato ad altri device non puo' ne ricevere e ne trasmettere frame.
Usa il concetto di enslaving: ossia le vengono assegnate delle
interfacce dette \textbf{enslaved} da controllare.

Quando si tenta di trasmettere su un bridge (es:
\passthrough{\lstinline!br0!}), il device driver gestisce i device
\textbf{enslaved} effettuando il forwarding dei frame.

Posso eventualmente nell'host, assegnare un'indirizzo IP al bridge
facendolo comportare come gateway per le interfacce collegate alle sue
porte. Dunque un brdige puo' essere esteso per partecipare al Layer 3.

\subsection[Netfilter (Capitolo 5 e 9 IPV4 routing)
]{\texorpdfstring{Netfilter (Capitolo 5 e 9 IPV4 routing) \footnote{\label{fn:ref_25}https://www.digitalocean.com/community/tutorials/a-deep-dive-into-iptables-and-netfilter-architecture}
\footnote{\label{fn:ref_26}https://linuxhandbook.com/iptables-vs-nftables/}}{Netfilter (Capitolo 5 e 9 IPV4 routing)  }}\label{netfilter-capitolo-5-e-9-ipv4-routing-3-4}

E' un sottosistema del kernel che fornisce le seguenti funzionalita
all'interno del sottosistema di routing:

\begin{itemize}
\tightlist
\item
  Selezione e filtraggio dei pacchetti
\item
  NAT
\item
  Tracciamento delle connessioni
\item
  Accounting di rete.
\end{itemize}

Per quanto riguarda il filtraggio dei pacchetti, abbiamo a disposizione
5 hook differenti:

\begin{itemize}
\tightlist
\item
  \passthrough{\lstinline!NF\_INET\_PRE\_ROUTING!}: e' chiamato nel
  metodo \passthrough{\lstinline!ip\_rcv()!} del kernel, quinid ancora
  prima di entrare nel sottosistema di routing.
\item
  \passthrough{\lstinline!NF\_INET\_LOCAL\_IN!}: e' chiamato nel momento
  in cui il pacchetto e' indirizzato al local host.
\item
  \passthrough{\lstinline!NF\_INET\_FORWARD!}: e' chiamato dopo che e'
  stato effettuato il lookup nel sottosistema di routing.
\item
  \passthrough{\lstinline!NF\_INET\_POST\_ROUTING!}: e' chiamato per i
  pacchetti che devono essere inoltrati, pochi istanti prima che il
  pacchetto entri fisicamente nel cavo.
\item
  \passthrough{\lstinline!NF\_INET\_LOCAL\_OUT!}: e' chiamato per i
  pacchetti creati nel local host, prima che escano.
\end{itemize}

Molte utility, comunemente usate, si basano sul sottosistema netfilter:

\begin{itemize}
\tightlist
\item
  \passthrough{\lstinline!iptables!}: agisce come front-end per
  netfilter ed implementa un \textbf{firewall}. Fornisce un'interfaccia
  per aggiungere, eliminare, regole netfilter, aggiungere tabelle,
  ecc\ldots{}
\item
  \passthrough{\lstinline!nftables!}: e' sviluppato da
  \textbf{netfilter} (l'organizzazione dietro al sottosistema
  netfilter). E' un rimpiazzo per iptables che e' riconosciuto avere
  problemi di scalabilita e performance.
\end{itemize}

\textbf{iptables} utilizza le \textbf{tables} per organizzare i filtri
da applicare, suddividendo le regole in classi di \textbf{table}. In
BlueAgent abbiamo utilizzato per esempio la classe
\passthrough{\lstinline!nat!}, usata per definire le regole per
applicare il NAT.

Spesso le connessioni a livello applicativo hanno una connessione
mantenuta nel tempo ed una serie di eventi che determinano la storia
della comunicazione. Il \textbf{Connection Tracking}
(\passthrough{\lstinline!cpnntrack!}) Layer permette al kernel di tenere
traccia delle sessioni create, e' la struttura dati fondamentale per
implementare il NAT. \passthrough{\lstinline!conntrack!} mantiene per
ogni connessione che deve usufruire del NAT le traduzioni necessarie per
gli indirizzi IP.
